\documentclass[11pt]{article}
\usepackage[margin=2.6cm]{geometry}
\usepackage{amsmath,amssymb,amsthm}
\usepackage{booktabs}
\usepackage{array}
\usepackage{microtype}
\usepackage{xcolor}
\usepackage[colorlinks=true,linkcolor=blue!50!black,citecolor=blue!50!black]{hyperref}

\newtheorem{theorem}{Theorem}[section]
\newtheorem{proposition}[theorem]{Proposition}
\theoremstyle{remark}
\newtheorem{remark}[theorem]{Remark}

\newcommand{\C}{\mathbb{C}}
\newcommand{\Q}{\mathbb{Q}}
\newcommand{\Z}{\mathbb{Z}}
\newcommand{\Nm}{\operatorname{Nm}}
\newcommand{\tr}{\operatorname{tr}}
\newcommand{\ind}{\operatorname{ind}}
\newcommand{\T}[1]{T_{#1}}
\newcommand{\rh}[1]{\rho(T_{#1})}
\newcommand{\mat}[4]{\begin{pmatrix} #1 & #2 \\ #3 & #4 \end{pmatrix}}

\title{Explicit irreducible representations of the Hecke algebras\\
of the cubic arc-transitive catalogue}
\author{Fable (Claude), with Robin Langer}
\date{June 2026}

\begin{document}
\maketitle

\begin{abstract}
For every $(G,H)$ coset pair and every graph in the catalogue of cubic
arc-transitive graphs (\texttt{CATALOGUE.md}), we compute all irreducible
representations of the Hecke algebra $\mathcal H(G,H)=\operatorname{End}_G \C[G/H]$
explicitly: spherical function tables for the Gelfand pairs, and exact matrix
entries (over cyclotomic fields) for every irreducible representation of
dimension $\ge 2$ in the non-Gelfand cases. The computation independently
re-derives the Meinhold-72 relations $A^2=I$, $B^2=13I$, $AB+BA=-4I$, and
shows that no universal formula $N=k+2\sqrt{k-\mu}$ governs the scalar $N$
in $B^2=NI$: the strongly-regular-graph mechanism behind $13=7+2\cdot 3$
is special to the $G_2(2)$ geometry. Companion GAP scripts:
\texttt{hecke\_irreps.g}, \texttt{hecke\_irreps2.g}, \texttt{hecke\_spherical.g}.
\end{abstract}

\tableofcontents

%=====================================================================
\section{Method}\label{sec:method}

\subsection{The Hecke algebra of a coset pair}

Let $G$ be a finite group, $H\le G$, $n=[G:H]$, and let
$\pi=\ind_H^G \mathbf 1$ be the permutation character of $G$ on $\C[G/H]$.
The \emph{Hecke algebra} is
\[
  \mathcal H(G,H) \;=\; \operatorname{End}_G \C[G/H]
  \;\cong\; \bigoplus_{\chi} M_{m_\chi}(\C),
  \qquad m_\chi = \langle \pi, \chi\rangle .
\]
A basis is given by the \emph{orbital matrices} $A_1,\dots,A_r$: the $0/1$
matrices of the orbits of $G$ on $(G/H)^2$, equivalently the adjacency
matrices of the suborbits of the point stabilizer. The rank satisfies
$r=\sum_\chi m_\chi^2$, and $(G,H)$ is a \emph{Gelfand pair} iff every
$m_\chi\le 1$, iff $\mathcal H$ is commutative.

By double-centralizer duality the irreducible representations of
$\mathcal H(G,H)$ are indexed by the constituents $\chi$ of $\pi$; the
Hecke irrep attached to $\chi$ has dimension $m_\chi$ (the multiplicity,
\emph{not} the degree $\chi(1)$), and occurs in $\C[G/H]$ with multiplicity
$\chi(1)$. For $m_\chi=1$ the irrep is the \emph{spherical function}
$\omega_\chi:\T i\mapsto \omega_\chi(\T i)\in\C$.

\subsection{Exact extraction of the $m$-dimensional irreps}

All matrices in this paper were obtained by the following method, which
stays exact over cyclotomic fields and avoids the MeatAxe.

\begin{enumerate}
\item Compute the isotypic projector
  $P_\chi = \frac{\chi(1)}{|G|}\sum_{g\in G}\chi(g^{-1})\,\pi(g)$
  via class-sum matrices. Its image $W$ is the $\chi$-isotypic component,
  $W \cong U_\chi \otimes \C^{\chi(1)}$ where $U_\chi$ is the
  $m$-dimensional Hecke irrep ($A_i$ acts as $\rho(\T i)\otimes 1$).
\item Find an eigenvector: pick an orbital matrix $A_a$ and an eigenvalue
  $\lambda$ of degree $\le 2$ over $\Q$ such that
  $\dim\bigl(\ker(A_a-\lambda)\cap W\bigr) = \chi(1)$. This dimension count
  forces $\lambda$ to be a \emph{simple} eigenvalue of $\rho(\T a)$ on
  $U_\chi$, so every vector of the intersection is a pure tensor
  $e_\lambda\otimes u$.
\item Spin: the closure of one such vector under $\{A_i\}$ is then a
  \emph{single} $m$-dimensional irreducible $\mathcal H$-module
  $U_\chi\otimes u$. Restricting each $A_i$ to it gives exact
  $m\times m$ matrices.
\item Normalize: change to an eigenbasis of one generator, and for $m=2$
  rescale so that some off-diagonal generator becomes symmetric.
\item Verify: the algebra generated by $\{\rho(\T i)\}$ has dimension
  $m^2$, and $\rho\bigl(\sum_i \T i\bigr)=0$ (the all-ones matrix dies on
  every nontrivial block). Both checks passed in every case below.
\end{enumerate}

\subsection{Conventions and notation}

$\T1$ is the identity orbital; suborbits are sorted by length (this
numbering differs from \texttt{MEINHOLD-HECKE.md} / \texttt{why13.tex};
the dictionary is given in \S\ref{sec:meinhold}). Matrices act on row
vectors. Recurring quadratic irrationalities:
\[
\begin{array}{llll}
\toprule
\text{symbol} & \text{value} & \text{satisfies} & \text{field}\\
\midrule
\alpha & \tfrac{-1+\sqrt{-7}}{2} = \zeta_7+\zeta_7^2+\zeta_7^4
       & \alpha^2+\alpha+2=0,\ \alpha\bar\alpha=2 & \Q(\sqrt{-7})\\[2pt]
\eta_0 & \tfrac{-1+\sqrt{13}}{2} = \textstyle\sum_{a\in\mathrm{QR}_{13}}\zeta_{13}^a
       & \eta^2+\eta-3=0,\ \eta_0\eta_1=-3 & \Q(\sqrt{13})\\[2pt]
\eta_1 & \tfrac{-1-\sqrt{13}}{2} & \eta_0-\eta_1=\sqrt{13} & \Q(\sqrt{13})\\[2pt]
\sqrt5 & (\zeta_5+\zeta_5^4)-(\zeta_5^2+\zeta_5^3) & \text{Gauss sum mod }5 & \Q(\sqrt5)\\
\bottomrule
\end{array}
\]

%=====================================================================
\section{The Gelfand backbone: spherical function tables}\label{sec:gelfand}

Every distance-transitive graph in the catalogue yields a Gelfand pair
$(\operatorname{Aut}\Gamma, \operatorname{Stab}(v))$ whose Hecke algebra is
the Bose--Mesner algebra of the distance scheme: rank $=$ diameter $+1$,
and the spherical values $\omega_\chi(\T i)$ on the distance-$i$ orbital
are the entries of the $P$-matrix (graph eigenvalue data). The tables below
were computed by \texttt{hecke\_spherical.g}; rows are spherical functions,
columns the orbitals $\T1,\dots,\T r$ sorted by suborbit length.

Below, each table lists the spherical functions $\omega_\chi$ (rows,
labelled by the constituent $\chi$ and its degree) against the orbitals
(columns, headed by suborbit lengths). Cyclotomic outputs have been
converted to surds; each conversion was verified in GAP
($E(8)-E(8)^3=\sqrt2$, $-E(12)^7+E(12)^{11}=\sqrt3$,
$-E(5)+E(5)^2+E(5)^3-E(5)^4=-\sqrt5$, etc.).

\subsection{Complete graphs and the cube family}

\textbf{$K_4$} (rank 2; suborbits $1,3$):\quad
$\chi_{\deg 1}: (1,3)$;\quad $\chi_{\deg 3}: (1,-1)$.

\medskip
\textbf{$K_{3,3}$} (rank 3; suborbits $1,2,3$):
\[
\begin{array}{l|rrr}
 & 1 & 2 & 3\\
\hline
\deg 1 & 1 & 2 & 3\\
\deg 1 & 1 & 2 & -3\\
\deg 4 & 1 & -1 & 0\\
\end{array}
\qquad
\textbf{Cube}\ (\text{rank }4):\quad
\begin{array}{l|rrrr}
 & 1 & 1 & 3 & 3\\
\hline
\deg 1 & 1 & 1 & 3 & 3\\
\deg 1 & 1 & -1 & -3 & 3\\
\deg 3 & 1 & 1 & -1 & -1\\
\deg 3 & 1 & -1 & 1 & -1\\
\end{array}
\]

\subsection{Petersen, Heawood, M\"obius--Kantor, Pappus, Desargues}

\[
\textbf{Petersen}:\;
\begin{array}{l|rrr}
 & 1 & 3 & 6\\
\hline
\deg 1 & 1 & 3 & 6\\
\deg 4 & 1 & -2 & 1\\
\deg 5 & 1 & 1 & -2\\
\end{array}
\qquad
\textbf{Heawood}:\;
\begin{array}{l|rrrr}
 & 1 & 3 & 4 & 6\\
\hline
\deg 1 & 1 & 3 & 4 & 6\\
\deg 1 & 1 & -3 & -4 & 6\\
\deg 6 & 1 & -\sqrt2 & \sqrt2 & -1\\
\deg 6 & 1 & \sqrt2 & -\sqrt2 & -1\\
\end{array}
\]
(The identical table holds for $\mathrm{PSL}(2,5)/S_3$ --- Petersen as a
coset graph --- and the Heawood data realizes the eigenvalues
$\pm\sqrt2$ of the incidence graph of $\mathrm{PG}(2,2)$.)

\[
\textbf{M\"obius--Kantor}:\;
\begin{array}{l|rrrrrr}
 & 1 & 1 & 2 & 3 & 3 & 6\\
\hline
\deg 1 & 1 & 1 & 2 & 3 & 3 & 6\\
\deg 1 & 1 & 1 & -2 & -3 & -3 & 6\\
\deg 3 & 1 & 1 & -2 & 1 & 1 & -2\\
\deg 3 & 1 & 1 & 2 & -1 & -1 & -2\\
\deg 4 & 1 & -1 & 0 & \sqrt3 & -\sqrt3 & 0\\
\deg 4 & 1 & -1 & 0 & -\sqrt3 & \sqrt3 & 0\\
\end{array}
\]
\[
\textbf{Pappus}:\;
\begin{array}{l|rrrrr}
 & 1 & 2 & 3 & 6 & 6\\
\hline
\deg 1 & 1 & 2 & 3 & 6 & 6\\
\deg 1 & 1 & 2 & -3 & 6 & -6\\
\deg 4 & 1 & 2 & 0 & -3 & 0\\
\deg 6 & 1 & -1 & \sqrt3 & 0 & -\sqrt3\\
\deg 6 & 1 & -1 & -\sqrt3 & 0 & \sqrt3\\
\end{array}
\]

\[
\textbf{Desargues}:\;
\begin{array}{l|rrrrrr}
 & 1 & 1 & 3 & 3 & 6 & 6\\
\hline
\deg 1 & 1 & 1 & 3 & 3 & 6 & 6\\
\deg 1 & 1 & -1 & -3 & 3 & 6 & -6\\
\deg 4 & 1 & 1 & -2 & -2 & 1 & 1\\
\deg 4 & 1 & -1 & 2 & -2 & 1 & -1\\
\deg 5 & 1 & 1 & 1 & 1 & -2 & -2\\
\deg 5 & 1 & -1 & -1 & 1 & -2 & 2\\
\end{array}
\]

\subsection{Dodecahedron and Nauru}

\[
\textbf{Dodecahedron}:\;
\begin{array}{l|rrrrrr}
 & 1 & 1 & 3 & 3 & 6 & 6\\
\hline
\deg 1 & 1 & 1 & 3 & 3 & 6 & 6\\
\deg 3 & 1 & -1 & -\sqrt5 & \sqrt5 & 2 & -2\\
\deg 3 & 1 & -1 & \sqrt5 & -\sqrt5 & 2 & -2\\
\deg 4 & 1 & 1 & -2 & -2 & 1 & 1\\
\deg 4 & 1 & -1 & 0 & 0 & -3 & 3\\
\deg 5 & 1 & 1 & 1 & 1 & -2 & -2\\
\end{array}
\]
\[
\textbf{Nauru}:\;
\begin{array}{l|rrrrrrr}
 & 1 & 2 & 3 & 3 & 3 & 6 & 6\\
\hline
\deg 1 & 1 & 2 & 3 & 3 & 3 & 6 & 6\\
\deg 1 & 1 & 2 & -3 & -3 & 3 & 6 & -6\\
\deg 3 & 1 & 2 & -1 & 3 & -1 & -2 & -2\\
\deg 3 & 1 & 2 & 1 & -3 & -1 & -2 & 2\\
\deg 4 & 1 & -1 & 0 & 0 & 3 & -3 & 0\\
\deg 6 & 1 & -1 & 2 & 0 & -1 & 1 & -2\\
\deg 6 & 1 & -1 & -2 & 0 & -1 & 1 & 2\\
\end{array}
\]
Note the dodecahedron's Aut-table (here, $\mathrm{Aut}=C_2\times A_5$,
rank 6) versus its rank-8 refinement under $\mathrm{PSL}(2,5)$ acting on
cosets of $C_3$ (\S\ref{sec:psl25}): the smaller group splits two
length-6 suborbits into four of length 3 and turns the pair Gelfand
$\to$ non-Gelfand.

\subsection{Tutte--Coxeter, the 12-cage, and $GH(2,2)$}

\[
\textbf{Tutte--Coxeter}:\;
\begin{array}{l|rrrrr}
 & 1 & 3 & 6 & 8 & 12\\
\hline
\deg 1 & 1 & 3 & 6 & 8 & 12\\
\deg 1 & 1 & -3 & 6 & 8 & -12\\
\deg 9 & 1 & -2 & 1 & -2 & 2\\
\deg 9 & 1 & 2 & 1 & -2 & -2\\
\deg 10 & 1 & 0 & -3 & 2 & 0\\
\end{array}
\]
\[
\textbf{12-cage orbits} = \textbf{$G_2(2)/P_1$}:\;
\begin{array}{l|rrrr}
 & 1 & 6 & 24 & 32\\
\hline
\deg 1 & 1 & 6 & 24 & 32\\
\deg 14 & 1 & -3 & 6 & -4\\
\deg 21 & 1 & 3 & 0 & -4\\
\deg 27 & 1 & -1 & -4 & 4\\
\end{array}
\]
The Tutte 12-cage is \emph{not} vertex-transitive (orbits $63+63$);
both orbits are the point/line actions of $G_2(2)$ on $GH(2,2)$ and
produce \emph{literally identical} spherical tables, even though no
automorphism exchanges them at $q=2$. The eigenvalue row
$(6,3^{21},-1^{27}\dots)$ matches \texttt{why13.tex} \S3.

\subsection{Coxeter-candidate coset pairs in $\mathrm{PGL}(2,7)$}

\[
\textbf{$\mathrm{PGL}(2,7)/A_4$} (n=28):\;
\begin{array}{l|rrrrrr}
 & 1 & 1 & 4 & 4 & 6 & 12\\
\hline
\deg 1 & 1 & 1 & 4 & 4 & 6 & 12\\
\deg 1 & 1 & 1 & -4 & -4 & -6 & 12\\
\deg 6 & 1 & 1 & -\sqrt2 & -\sqrt2 & 2\sqrt2 & -2\\
\deg 6 & 1 & 1 & \sqrt2 & \sqrt2 & -2\sqrt2 & -2\\
\deg 7 & 1 & -1 & -2 & 2 & 0 & 0\\
\deg 7 & 1 & -1 & 2 & -2 & 0 & 0\\
\end{array}
\]
\[
\textbf{$\mathrm{PGL}(2,7)/D_{12}$} (n=28):\;
\begin{array}{l|rrrrr}
 & 1 & 3 & 6 & 6 & 12\\
\hline
\deg 1 & 1 & 3 & 6 & 6 & 12\\
\deg 6 & 1 & -1+\sqrt2 & 2+\sqrt2 & -2\sqrt2 & -2\\
\deg 6 & 1 & -1-\sqrt2 & 2-\sqrt2 & 2\sqrt2 & -2\\
\deg 7 & 1 & -1 & -2 & -2 & 4\\
\deg 8 & 1 & 2 & -2 & 1 & -2\\
\end{array}
\]

\subsection{Biggs--Smith: $\mathrm{PSL}(2,17)/S_4$, $n=102$}

Rank 8, suborbits $(1,3,6,8,12,24,24,24)$. Three Galois orbits of
spherical functions: the trivial one; a $\Q(\sqrt{17})$ pair (deg 9);
a cubic triple (deg 16) over $\Q(\zeta_9)^+$; and two rational rows
(deg 17, 18).
\[
\begin{array}{l|rrrrrrrr}
 & 1 & 3 & 6 & 8 & 12 & 24 & 24 & 24\\
\hline
\deg 1 & 1 & 3 & 6 & 8 & 12 & 24 & 24 & 24\\
\deg 9 & 1 & \frac{1+\sqrt{17}}{2} & \frac{3+\sqrt{17}}{2} & 1-\sqrt{17} & 4 & -1+\sqrt{17} & -8 & 1-\sqrt{17}\\[2pt]
\deg 9 & 1 & \frac{1-\sqrt{17}}{2} & \frac{3-\sqrt{17}}{2} & 1+\sqrt{17} & 4 & -1-\sqrt{17} & -8 & 1+\sqrt{17}\\[2pt]
\deg 16 & 1 & -\beta-1 & \beta^2+2\beta-2 & \beta^2+\beta-4 & -3\beta^2-\beta+5 & 2\beta^2+\beta-4 & -2\beta^2+5 & \beta^2-2\beta\\
\deg 17 & 1 & 0 & -3 & 2 & 0 & 6 & 0 & -6\\
\deg 18 & 1 & 2 & 1 & 2 & -2 & -6 & 4 & -2\\
\end{array}
\]
where $\beta=\zeta_9+\zeta_9^{-1}$ runs over the three roots of
$\beta^3-3\beta+1=0$ (one row per root; the orbital order of the three
length-24 columns is permuted between conjugate rows --- see
\texttt{hecke\_spherical.out} for the exact alignment). Substituting
$\beta=-(t+1)$ recovers the Biggs--Smith adjacency cubic $t^3+3t^2-3$:
the spectrum
$\{3,\ \tfrac{1\pm\sqrt{17}}{2}^{\,(9)},\ (t^3+3t^2-3)^{(16\times3)},\
0^{(17)},\ 2^{(18)}\}$ is reassembled block by block.

\subsection{Two-transitive geometries (rank 2) and $GQ(3,3)$}

\[
\begin{array}{l|cc}
\text{pair} & \omega_{\text{triv}} & \omega_{\text{Steinberg-like}}\\
\hline
\mathrm{PSL}(3,3)/P_1\ (13 \text{ pts } \mathrm{PG}(2,3)) & (1,12) & (1,-1)\\
\mathrm{PSL}(3,4)/P_1\ (21 \text{ pts } \mathrm{PG}(2,4)) & (1,20) & (1,-1)\\
\mathrm{PSU}(3,3)/H_{28}\ (\text{Hermitian unital}) & (1,27) & (1,-1)\\
\mathrm{Sz}(8)/H_{65}\ (\text{Suzuki ovoid}) & (1,64) & (1,-1)\\
\end{array}
\]
\[
\textbf{$\mathrm{PSp}(4,3)/P_1$} (40 \text{ pts of } GQ(3,3)):\;
\begin{array}{l|rrr}
 & 1 & 12 & 27\\
\hline
\deg 1 & 1 & 12 & 27\\
\deg 15 & 1 & -4 & 3\\
\deg 24 & 1 & 2 & -3\\
\end{array}
\]

%=====================================================================
\section{$\mathrm{PSL}(2,5)/C_3$, $n=20$: the dodecahedral miniature}
\label{sec:psl25}

$\mathcal H \cong \C^4\oplus M_2(\C)$, rank $8$, suborbits
$(1,1,3,3,3,3,3,3)$. The coset graph of the valency-3 orbital $\T3$ is
the \emph{dodecahedron}.

\subsection{Spherical functions}

\[
\begin{array}{l|rrrrrrrr}
 & \T1 & \T2 & \T3 & \T4 & \T5 & \T6 & \T7 & \T8\\
\hline
\chi_1\ (\deg 1) & 1 & 1 & 3 & 3 & 3 & 3 & 3 & 3\\
\chi_2\ (\deg 3) & 1 & -1 & -\sqrt5 & \sqrt5 & 1 & -1 & 1 & -1\\
\chi_3\ (\deg 3) & 1 & -1 & \sqrt5 & -\sqrt5 & 1 & -1 & 1 & -1\\
\chi_5\ (\deg 5) & 1 & 1 & 1 & 1 & -1 & -1 & -1 & -1\\
\end{array}
\]

\subsection{The 2-dimensional irrep (on $\chi_4$, $\deg 4$)}

In the eigenbasis of $\rho(\T2)$:
\[
\rh1 = I,\quad
\rh2 = \mat{1}{0}{0}{-1},\quad
\rh3 = \rh4 = \mat{-2}{0}{0}{0},
\]
\[
\rh5 = \mat{1/2}{1}{-15/4}{-3/2},\quad
\rh6 = \mat{1/2}{-1}{-15/4}{3/2},\quad
\rh7 = \mat{1/2}{-1}{15/4}{-3/2},\quad
\rh8 = \mat{1/2}{1}{15/4}{3/2}.
\]
Characteristic polynomials: $\rh2$: $x^2-1$; $\rh3=\rh4$: $x(x+2)$;
$\rh5,\rh7$: $x^2+x+3$ (complex, $\det = 3 =$ valency);
$\rh6,\rh8$: $(x-3)(x+1)$.

\begin{proposition}[Relations]\label{prop:psl25}
Set $A=\rh2$, $C=\rh5$. Then
\[
A^2 = I,\qquad C^2+C+3I = 0,\qquad
AC = \rh8,\quad CA = \rh6,\quad ACA = \rh7,
\]
and the closest analogue of $B^2=13I$ is the shifted square
\[
\bigl(\rh6 - I\bigr)^2 = 4I .
\]
\end{proposition}

\begin{remark}
The block does not separate $\T3$ from $\T4$ ($\rh3=\rh4$); they are
separated by $\chi_2,\chi_3$ via $\pm\sqrt5$. Reassembling the spectrum of
the symmetric orbital $\T3$ from all blocks gives
$\{3 \mid \pm\sqrt5 \mid 0,-2 \mid 1\}$ with multiplicities
$\{1\mid 3,3\mid 4,4\mid 5\}$ --- exactly the dodecahedron spectrum
$3,\ \sqrt5^{\,3},\ 1^5,\ 0^4,\ (-2)^4,\ (-\sqrt5)^{3}$. The famous
$\pm\sqrt5$ of the dodecahedron lives entirely in the
\emph{one-dimensional} Hecke irreps (the degree-3 characters of $A_5$);
the $2\times2$ block contributes only integer eigenvalues.
\end{remark}

%=====================================================================
\section{$\mathrm{PSL}(2,7)/C_3$, $n=56$}\label{sec:psl27}

$\mathcal H \cong \C^3\oplus M_2\oplus M_2\oplus M_3$, rank $20$,
suborbits $(1,1,3^{18})$.

\subsection{The $\Q(\sqrt{-7})$ block (on $\chi_4$, $\deg 6$)}

Here $\rho(\T2)=I$: the free involution acts trivially on this block.
With $\alpha=\frac{-1+\sqrt{-7}}{2}$, the distinct images are
\[
\begin{array}{lll}
\toprule
\text{matrix} & \text{which } \T i & \text{charpoly}\\
\midrule
\operatorname{diag}(\alpha,\bar\alpha) & \T3,\T5,\T{12},\T{17} & x^2+x+2\\
\operatorname{diag}(\bar\alpha,\alpha) & \T4,\T6,\T7,\T{14} & x^2+x+2\\
\mat{1}{2}{2}{1} & \T8,\T{18} & (x-3)(x+1)\\
\mat{1}{\alpha-1}{\bar\alpha-1}{1} & \T9,\T{15} & (x-3)(x+1)\\
\mat{0}{\bar\alpha}{\alpha}{0} & \T{10},\T{11},\T{16},\T{20} & x^2-2\\
\mat{-1}{-\bar\alpha}{-\alpha}{-1} & \T{13},\T{19} & x^2+2x-1\\
\bottomrule
\end{array}
\]

\begin{proposition}\label{prop:psl27a}
Set $D=\operatorname{diag}(\alpha,\bar\alpha)$ and
$S=\bigl(\begin{smallmatrix}0&\bar\alpha\\ \alpha&0\end{smallmatrix}\bigr)$.
Then
\[
D^2+D+2I=0,\qquad \boxed{S^2 = 2I},\qquad SDS^{-1}=\bar D,
\qquad \rho(\T{13}) = -I - S,
\]
where $2=\alpha\bar\alpha=\Nm(\alpha)$ is the prime $2$, split in
$\Q(\sqrt{-7})$. The six matrices above span $M_2(\C)$, so $\{D,S\}$
generate the block. The eigenvalues $-1\pm\sqrt2$ of $\rho(\T{13})$
exhibit $\sqrt2$ as the ``junior $\sqrt{13}$''.
\end{proposition}

\subsection{The integral $3$-dimensional block (on $\chi_5$, $\deg 7$)}

In the eigenbasis of $\rho(\T2)=\operatorname{diag}(1,-1,-1)$ \emph{all
twenty} images are integer matrices. Representative images:
\[
\rh3 = \begin{pmatrix} 1&1&0\\ 2&0&2\\ 4&0&1\end{pmatrix},\quad
\rh4 = \begin{pmatrix} 1&0&1\\ 4&1&0\\ 2&2&0\end{pmatrix},\quad
\rh8 = \begin{pmatrix} -1&0&0\\ 0&1&0\\ 0&2&-3\end{pmatrix},\quad
\rh{10} = \begin{pmatrix} -1&-1&1\\ 2&0&0\\ -2&0&0\end{pmatrix},
\]
\[
\rh{11} = \begin{pmatrix} -1&1&-1\\ 2&0&0\\ -2&0&0\end{pmatrix},\quad
\rh{13} = \begin{pmatrix} -1&0&0\\ 0&-1&2\\ 0&-2&1\end{pmatrix},
\]
the remaining ones being sign/transpose variants (full list in
\texttt{hecke\_irreps.out}). Characteristic polynomials:
\[
\begin{array}{ll}
\toprule
\text{charpoly} & \text{which }\T i\\
\midrule
x^3-2x^2-x-6 & \T3,\T4,\T{12},\T{14}\\
x^3+x-6\ \text{(irreducible over }\Q\text{: genuinely cubic eigenvalues)} & \T5,\T6,\T7,\T{17}\\
(x+3)(x-1)(x+1) & \T8,\T{15}\\
(x-3)(x+1)^2 & \T9,\T{18}\\
x(x^2+x+4) & \T{10},\T{16}\\
x(x^2+x-4)\quad\text{eigenvalues } 0,\ \tfrac{-1\pm\sqrt{17}}{2} & \T{11},\T{20}\\
(x+1)(x^2+3)\quad\text{eigenvalues } -1,\ \pm i\sqrt3 & \T{13},\T{19}\\
\bottomrule
\end{array}
\]

\subsection{The rational block (on $\chi_6$, $\deg 8$)}

In the eigenbasis of $\rho(\T2)=\operatorname{diag}(1,-1)$:
\[
\begin{array}{lll}
\toprule
\text{matrix} & \text{which }\T i & \text{charpoly}\\
\midrule
\mat{-1/2}{1}{-7/4}{-1/2} & \T3,\T{14} & x^2+x+2\\[2pt]
\mat{-1/2}{-1}{7/4}{-1/2} & \T4,\T{12} & x^2+x+2\\[2pt]
\mat{-1/2}{1}{7/4}{1/2} & \T5,\T7 & x^2-2\\[2pt]
\mat{-1/2}{-1}{-7/4}{1/2} & \T6,\T{17} & x^2-2\\[2pt]
\operatorname{diag}(-1,1) & \T8,\T{15} & x^2-1\\
-I & \T9,\T{18} & (x+1)^2\\[2pt]
\mat{1/2}{1}{-7/4}{-3/2} & \T{10}\ (\T{16}\ \text{variant}) & x^2+x+1\quad(\rho^3=I)\\[2pt]
\mat{1/2}{-1}{-7/4}{3/2} & \T{11}\ (\T{20}\ \text{variant}) & x^2-2x-1\quad(1\pm\sqrt2)\\[2pt]
\mat{2}{0}{0}{0} & \T{13},\T{19} & x(x-2)\\
\bottomrule
\end{array}
\]

\begin{proposition}\label{prop:psl27c}
With $B=\rh5$: $\boxed{B^2=2I}$ --- the same norm-$2$ phenomenon as in
$\chi_4$, now realized over $\Q$. Moreover $\rho(\T{10})$ has
\emph{finite order $3$} (charpoly $x^2+x+1$): a Hecke generator acting as
a rotation. The eigenvalues $1\pm\sqrt2$ recur: the $\Z[\sqrt2]$ thread
runs through both $2$-dimensional blocks.
\end{proposition}

%=====================================================================
\section{F26A, $n=26$: Galois-conjugate blocks over $\Q(\sqrt{13})$}
\label{sec:f26a}

$\operatorname{Aut}=C_{13}\rtimes C_6$ (Frobenius, sharply arc-transitive),
$\mathcal H\cong\C^2\oplus M_2\oplus M_2$, rank $10$, suborbits
$(1,1,3^8)$.

\subsection{Spherical functions}
$\chi_1$: $(1,1,3,3,3,3,3,3,3,3)$;\quad
$\chi_2$: $(1,-1,-3,3,-3,3,-3,3,-3,3)$.

\subsection{The 2-dimensional irreps}

The two $M_2$ blocks (on the two degree-6 characters of
$C_{13}\rtimes C_6$) are \emph{Galois conjugate over} $\Q(\sqrt{13})$:
swapping $\eta_0\leftrightarrow\eta_1$ passes between them. On $\chi_7$,
in the eigenbasis of $\rho(\T2)=\operatorname{diag}(1,-1)$ and with
$\gamma := 2\eta_0+\tfrac54\eta_1$:
\[
\rh3 = \mat{\eta_0/2}{1}{-\gamma}{-\eta_0/2},\quad
\rh4 = \mat{\eta_1/2}{-(\eta_0+2\eta_1)}{-\sqrt{13}/4}{\eta_1/2},\quad
\rh5 = \mat{\eta_0/2}{-1}{\gamma}{-\eta_0/2},
\]
\[
\rh6 = \mat{\eta_0/2}{-1}{-\gamma}{\eta_0/2},\quad
\rh7 = \mat{\eta_1/2}{\eta_0+2\eta_1}{-\sqrt{13}/4}{-\eta_1/2},\quad
\rh8 = \mat{\eta_1/2}{\eta_0+2\eta_1}{\sqrt{13}/4}{\eta_1/2},
\]
\[
\rh9 = \mat{\eta_1/2}{-(\eta_0+2\eta_1)}{\sqrt{13}/4}{-\eta_1/2},\quad
\rh{10} = \mat{\eta_0/2}{1}{\gamma}{\eta_0/2}.
\]

\begin{proposition}[F26A relations]\label{prop:f26a}
With $2-\eta_0=\frac{5-\sqrt{13}}{2}$ and $2-\eta_1=\frac{5+\sqrt{13}}{2}$:
\[
\begin{array}{ll}
\toprule
\T i & \text{relation}\\
\midrule
\T3,\T5 & T^2 = \tfrac{5-\sqrt{13}}{2}\, I\\[2pt]
\T7,\T9 & T^2 = \tfrac{5+\sqrt{13}}{2}\, I\\[2pt]
\T4,\T8 & T^2 = \eta_1 T - (2-\eta_1) I\\[2pt]
\T6,\T{10} & T^2 = \eta_0 T - (2-\eta_0) I\\
\bottomrule
\end{array}
\]
The traceless generators square to the scalars
$N_\pm = \frac{5\mp\sqrt{13}}{2}$ with
\[
N_+N_- = 3\ (\text{the valency}),\qquad N_+ + N_- = 5 .
\]
\end{proposition}

\begin{remark}
This $13$ is the \emph{same numeral} as Meinhold's but enters by a
different door: here $26 = 2\cdot 13$ and $\sqrt{13}$ arises from the
quadratic Gauss periods $\eta_0,\eta_1$ of $\Q(\zeta_{13})$, whereas in
\S\ref{sec:meinhold} it arises from the strongly regular graph identity
$13 = 7+2\cdot3$. The two Galois-conjugate blocks share the rational
invariants $(3,5)$.
\end{remark}

%=====================================================================
\section{$G_2(2)/\mathrm{PSL}(3,2)$ (Meinhold-72): cross-check}
\label{sec:meinhold}

$\mathcal H\cong\C^2\oplus M_2\oplus\C^2$, rank $8$, suborbits
$(1,1,7,7,7,7,21,21)$. Numbering dictionary to
\texttt{why13.tex}/\texttt{MEINHOLD-HECKE.md}: here $\T2$ is the free
involution ``$A$''; $\T4/\T6$ is the symmetric valency-7 pair ``$B/B'$''
(old $\T3/\T5$); $\T3/\T5$ is the directed pair (old $\T4/\T6$).

\subsection{Spherical functions}
\[
\begin{array}{l|rrrrrrrr}
 & \T1 & \T2 & \T3 & \T4 & \T5 & \T6 & \T7 & \T8\\
\hline
\chi_1 & 1 & 1 & 7 & 7 & 7 & 7 & 21 & 21\\
\chi_2 & 1 & -1 & 7 & -7 & 7 & -7 & -21 & 21\\
\chi_{21} & 1 & -1 & 1 & -1 & 1 & -1 & 3 & -3\\
\chi_{21}' & 1 & 1 & 1 & 1 & 1 & 1 & -3 & -3\\
\end{array}
\]

\subsection{The 2-dimensional irrep (on $\chi_{14}$)}

In the eigenbasis of $\rho(\T2)$:
\[
\rh2 = \mat{1}{0}{0}{-1},\quad
\rh3 = \mat{-2}{1}{-9}{-2},\quad
\rh4 = \mat{-2}{-1}{-9}{2},\quad
\rh5 = \mat{-2}{-1}{9}{-2},
\]
\[
\rh6 = \mat{-2}{1}{9}{2},\quad
\rh7 = \mat{3}{0}{0}{-3},\quad
\rh8 = 3I .
\]
Charpolys: $\rh3,\rh5$: $x^2+4x+13$ (the $-2\pm3i$ pair,
$|-2+3i|^2=13$); $\rh4,\rh6$: $x^2-13$; $\rh7$: $x^2-9$; $\rh8$ scalar.

\begin{theorem}[Independent re-derivation of the why-13 algebra]
\label{thm:meinhold}
With $A=\rh2$ and $B=\rh4$:
\[
A^2 = I,\qquad \boxed{B^2 = 13I},\qquad AB+BA = -4I,
\]
and the directed pair is $BA = \rh3$, $AB = \rh5$. This reproduces
\textup{why13.tex} Theorem~8.1 by isotypic projection, a method sharing
nothing with the original orbital path-counting derivation.
\end{theorem}

%=====================================================================
\section{$\mathrm{PSL}(2,13)/S_3$, $n=182$}\label{sec:psl213}

$\mathcal H \cong \C^4 \oplus M_2^{\oplus 3} \oplus M_3^{\oplus 2}$,
rank $34 = 4 + 3\cdot4 + 2\cdot9$, suborbits $(1,1,3^4,6^{28})$.
All matrices below from \texttt{hecke\_irreps2.out} (both $S_3$-classes
give equivalent data; we use class \#1). Set
$c_k = \zeta_7^k+\zeta_7^{-k}$, so $c_1+c_2+c_3=-1$ and each $c_k$
satisfies $x^3+x^2-2x-1=0$.

\subsection{One-dimensional irreps}

Besides the trivial and sign rows, the two degree-7 constituents give
spherical functions with values in $\Q(\sqrt{13})$ via the \emph{same}
Gauss periods $\eta_0,\eta_1$ as F26A (\S\ref{sec:f26a}): e.g.\
$\omega_{\chi_2}(\T3) = -\eta_0 = \frac{1-\sqrt{13}}{2}$ and length-6
orbitals take values such as $2-\eta_1=\frac{5+\sqrt{13}}{2}$
(full vectors in \texttt{hecke\_irreps2.out}). Both $\sqrt{13}$'s ---
F26A's and this one --- come from $\Q(\zeta_{13})$, unlike Meinhold's.

\subsection{The three Galois-conjugate $M_2$ blocks (on $\chi_4,\chi_5,\chi_6$, $\deg 12$)}

The three degree-12 characters of $\mathrm{PSL}(2,13)$ are conjugate
under $\mathrm{Gal}(\Q(\zeta_7)^+/\Q)\cong C_3$, and so are the three
Hecke blocks: the substitution $c_k\mapsto c_{2k}$ maps the $\chi_4$
block to the $\chi_6$ block to the $\chi_5$ block. (Verified on the
charpolys of the valency-3 generator:
$x^2-(c_2{+}c_3)x-(c_1{+}2c_2{+}3c_3)$ for $\chi_4$,
$x^2-(c_1{+}c_3)x-(3c_1{+}c_2{+}2c_3)$ for $\chi_6$,
$x^2-(c_1{+}c_2)x-(2c_1{+}3c_2{+}c_3)$ for $\chi_5$.)
It therefore suffices to give the $\chi_4$ block; in the eigenbasis of
$\rho(\T2)$,
\[
\rh2 = \mat{1}{0}{0}{-1},\qquad
\rh3 = \begin{pmatrix}
\tfrac12 c_1+\tfrac12 c_2+c_3 & 1\\[2pt]
\tfrac32 c_1+\tfrac74 c_2+\tfrac{11}4 c_3 & -\tfrac12 c_1+\tfrac12 c_2
\end{pmatrix},
\]
and $\{\rho(\T2),\rho(\T3)\}$ already generate the block ($\rho(\T3)$ is
non-diagonal with distinct eigenvalues). The remaining 31 images are in
\texttt{hecke\_irreps2.out} (lines 60--307). Unlike all previous $M_2$
blocks, no orbital here squares to a rational scalar: the entries
genuinely need the cubic field $\Q(\zeta_7)^+$.

\subsection{The integral Steinberg block (on $\chi_7$, $\deg 13$)}

All 34 images are \emph{integer} matrices. In the eigenbasis of
$\rho(\T2)=\operatorname{diag}(1,1,-1)$:
\[
\rh3 = \begin{pmatrix} 2&-1&-2\\ 2&-2&-3\\ 1&0&0\end{pmatrix},\quad
\rh4 = \begin{pmatrix} 2&-1&2\\ 2&-2&3\\ 1&0&0\end{pmatrix},\quad
\rh7 = \begin{pmatrix} 1&0&3\\ 3&-1&8\\ 0&1&3\end{pmatrix},\quad
\rh{21} = \begin{pmatrix} -2&1&-3\\ -5&4&-1\\ 3&-1&-2\end{pmatrix}.
\]
Characteristic polynomial highlights:
\[
\begin{array}{ll}
\toprule
\text{charpoly} & \text{which }\T i\\
\midrule
x^3+1\quad\Longrightarrow\ \boxed{\rho(\T3)^3=-I}\ \text{(finite order 6!)} & \T3,\T5\\
x^3-4x-1\ \text{(irreducible, totally real)} & \T4,\T6\\
x^3+x-26\quad (26=2\cdot13=n/7) & \T{21},\T{26}\\
x^3-3x^2-9x+2,\ x^3+3x^2+7x-2,\ \dots & \text{(rest: see \texttt{.out})}\\
\bottomrule
\end{array}
\]
Note the echo of $\mathrm{PSL}(2,7)/C_3$'s integral $3$-dim block,
whose analogous cubic was $x^3+x-6$; here it is $x^3+x-26$.

\subsection{The $\deg 14$ block: a cube root of $5$}

On $\chi_9$ ($\deg 14$), in the eigenbasis of
$\rho(\T2)=\operatorname{diag}(1,1,-1)$:
\[
\rh3 = \begin{pmatrix} \tfrac12&\tfrac12&-\tfrac{11}4\\ 2&1&3\\ 1&0&-\tfrac32\end{pmatrix},
\qquad
\text{charpoly } x^3 - 5,\qquad \boxed{\rho(\T3)^3 = 5\,I}.
\]
$\rho(\T5)$ satisfies the same relation; $\rho(\T4),\rho(\T6)$ have
charpoly $x^3-3x^2-x+5$. This is the first \emph{cubic} analogue of
$B^2=13I$: a traceless valency-3 orbital whose \emph{cube} is the scalar
$5$. The eigenvalues are $\sqrt[3]{5}\,\zeta_3^j$ --- the Hecke algebra
of a cubic graph manufacturing a cube root of $5$.

%=====================================================================
\section{Omitted large cases}\label{sec:omitted}

Two non-Gelfand pairs of the census are omitted here:
$\mathrm{PSL}(2,31)/S_4$ ($n=620$, rank 33, six $M_2$ blocks --- the
first failure of Gelfand-ness in the $S_4$ tower; its $1$-dimensional
spherical values live in $\Q(\sqrt2)$) and $\mathrm{PSL}(2,19)/S_3$
($n=570$, rank 101, three $M_4$ blocks). The extraction machinery is in
place (\texttt{hecke\_irreps3.g}, which restricts the orbital matrices to
the $\chi$-isotypic subspace before the eigenvalue search); these can be
run when wanted.

%=====================================================================
\section{The pattern: which $N$ appears in $B^2=NI$?}\label{sec:pattern}

\[
\begin{array}{llll}
\toprule
\text{algebra} & \text{traceless generator} & B^2=N\,I & \text{where } N \text{ lives}\\
\midrule
G_2(2)/L_2(7),\ \chi_{14} & \text{valency-7 symmetric} & N=13=7+2\cdot3
  & \Z;\ |-2\pm3i|^2 = 13 \text{ in } \Z[i]\\
\mathrm{PSL}(2,7)/C_3,\ \chi_4,\chi_6 & \text{valency-3 (two blocks)} & N=2
  & \Z;\ 2=\Nm\bigl(\tfrac{-1+\sqrt{-7}}{2}\bigr)\\
\mathrm{F26A},\ \chi_7/\chi_8 & \text{valency-3} & N=\tfrac{5\mp\sqrt{13}}{2}
  & \Z[\eta];\ \Nm(N)=3=\text{valency}\\
\mathrm{PSL}(2,5)/C_3,\ \chi_4 & (\rho(\T6)-I)\ \text{shifted} & N=4
  & \Z;\ \text{complex pair has } \det=3\\
\bottomrule
\end{array}
\]

In every case the directed/complex orbitals have determinant or norm tied
to the valency, while the symmetric generator squares to an $N$ measuring
the interference between the two chirality classes; for Meinhold-72 the
precise statement is $NN^{\mathsf T} = 7I + 2A_{\mathrm{SRG}(36,21,12,12)}$,
hence $13 = 7+2\cdot3$.

\begin{remark}[No universal formula]
There is no formula $N = k + 2\sqrt{k-\mu}$ valid across cases:
$\mathrm{PSL}(2,7)/C_3$ has $N = 2 < k = 3$, so the strongly-regular-graph
mechanism is special to the Meinhold geometry. The general synthetic
meaning of $N$ remains open (\textup{why13.tex} \S9).
\end{remark}

Side findings worth keeping: a Hecke generator of finite order $3$
($\mathrm{PSL}(2,7)/C_3$, $\chi_6$); genuinely cubic eigenvalues
($x^3+x-6$) in the $3$-dimensional irrep; $\sqrt{17}$ in the same block;
the dodecahedron's $\sqrt5$ living entirely in the $1$-dimensional irreps.

\end{document}
